\documentclass[11pt]{article}
\usepackage[a4paper, margin=2cm]{geometry}
\usepackage{graphicx}
\usepackage{amsmath, amssymb}
\usepackage{algorithm}
\usepackage{algpseudocode}
\usepackage{booktabs}
\usepackage{enumitem}
\usepackage{hyperref}

\title{Interim Report: Robot Arm Simulation with Inverse Kinematics and 3D Projection}
\author{Roman Prokhorov, Nazar Pasichnyk, Mykola Balyk}
\date{March 2026}

\begin{document}

\maketitle

\section{Research Area}

Robotic manipulation is a central problem in robotics, computer animation, and industrial automation. A fundamental challenge in this field is \textbf{inverse kinematics} (IK): given a desired position of a robot's end-effector (e.g., a gripper), determine the joint angles that achieve it~\cite{mathworks}. Unlike the straightforward \textit{forward kinematics}---where joint angles directly yield the end-effector position via matrix multiplication---inverse kinematics requires solving a nonlinear system that may have zero, one, or infinitely many solutions.

This problem is deeply rooted in \textbf{linear algebra}. The kinematic chain of a robot arm is modeled as a product of $4 \times 4$ transformation matrices (rotations and translations in homogeneous coordinates). The iterative IK solution relies on the \textbf{Jacobian matrix}, its pseudoinverse, and \textbf{Singular Value Decomposition} (SVD) for numerical stability. Visualizing the arm in real time additionally requires perspective projection matrices and orthonormal basis construction for the camera---all core linear algebra operations.

This project is interesting because it connects abstract algebraic concepts (matrix groups, pseudoinverses, SVD) directly to a tangible, visual application: a moving robot arm that tracks a target in 3D space.

\section{Aim and Tasks}

The aim is to develop an \textbf{interactive 3D robot arm simulator with inverse kinematics} that visually demonstrates how linear algebra techniques enable a robotic arm to reach user-specified target positions in real time. The system will combine forward kinematics (matrix chain multiplication), inverse kinematics (Jacobian pseudoinverse with SVD), and 3D-to-2D projection (perspective transformation) into a single cohesive application where users can drag a target point and watch the arm adjust its joint angles to track it.

\section{Literature Review and Chosen Method}

Several \textit{complete} approaches exist for solving inverse kinematics. Each represents a self-contained method to compute joint angles from a target position:

\begin{itemize}[nosep]
    \item \textbf{Analytical (closed-form) solutions} exist for specific arm geometries (e.g., 6-DOF arms with spherical wrists)~\cite{roboticworx}. They are fast but not generalizable.
    \item \textbf{Jacobian Transpose method}: updates angles via $\Delta\theta = \alpha \, J^T \Delta p$. Simple and cheap, but convergence is slow and depends on step size~\cite{buss2004}.
    \item \textbf{Jacobian Pseudoinverse method}: $\Delta\theta = J^+ \Delta p$ where $J^+ = J^T(JJ^T)^{-1}$. Converges faster and provides a least-norm solution, but is unstable near singularities~\cite{buss2004}.
    \item \textbf{Damped Least Squares (DLS)}: $\Delta\theta = J^T(JJ^T + \lambda^2 I)^{-1}\Delta p$. Adds a damping factor $\lambda$ to avoid blow-up at singularities~\cite{buss2004}.
    \item \textbf{CCD (Cyclic Coordinate Descent)}: a heuristic iterative method that adjusts one joint at a time. Easy to implement, but can produce unnatural motion~\cite{buss2004}.
    \item \textbf{FABRIK}: a fast geometric method that operates on positions rather than angles~\cite{buss2004}.
\end{itemize}

We choose the \textbf{Damped Least Squares} approach with \textbf{SVD-based pseudoinverse} computation. This method directly showcases key linear algebra concepts (Jacobian, SVD, matrix inversion, transpose), converges reliably, and handles singularities gracefully by thresholding small singular values:
\[
J = U \Sigma V^T, \qquad J^+ = V \Sigma^+ U^T,
\]
where $\Sigma^+$ replaces each singular value $\sigma_i$ with $\sigma_i / (\sigma_i^2 + \lambda^2)$.

For 3D rendering, we use the standard perspective projection pipeline with a $4\times 4$ projection matrix followed by perspective divide.

\section{Pseudocode}

\begin{algorithm}[H]
\caption{Inverse Kinematics via Damped Least Squares}
\begin{algorithmic}[1]
\Require Joint angles $\theta = [\theta_1, \dots, \theta_n]$, target position $p_{\text{target}}$, damping $\lambda$, tolerance $\varepsilon$
\Ensure Updated joint angles $\theta$
\While{$\| p_{\text{end}} - p_{\text{target}} \| > \varepsilon$}
    \State $p_{\text{end}} \gets \text{ForwardKinematics}(\theta)$ \Comment{chain of $4\times4$ matrix multiplications}
    \State $\Delta p \gets p_{\text{target}} - p_{\text{end}}$
    \State Compute Jacobian $J$: column $j = \text{axis}_j \times (p_{\text{end}} - p_j)$
    \State $U, \Sigma, V^T \gets \text{SVD}(J)$ \Comment{decompose the $3 \times n$ Jacobian}
    \State $\Sigma^+ \gets \text{diag}\!\left(\frac{\sigma_i}{\sigma_i^2 + \lambda^2}\right)$ \Comment{damped inverse of singular values}
    \State $J^+ \gets V \, \Sigma^+ U^T$
    \State $\Delta\theta \gets J^+ \cdot \Delta p$
    \State $\theta \gets \theta + \Delta\theta$
\EndWhile
\end{algorithmic}
\end{algorithm}

\begin{algorithm}[H]
\caption{Rendering Pipeline (per frame)}
\begin{algorithmic}[1]
\State Compute all joint world positions $\{p_0, p_1, \dots, p_n\}$ via FK (chain of $4\times4$ transforms)
\State \textbf{Construct View matrix} $V$:
\State \quad $\vec{z} \gets \text{normalize}(\text{camera\_pos} - \text{look\_at})$ \Comment{camera forward direction}
\State \quad $\vec{x} \gets \text{normalize}(\text{up} \times \vec{z})$ \Comment{camera right direction}
\State \quad $\vec{y} \gets \vec{z} \times \vec{x}$ \Comment{camera up direction}
\State \quad $V \gets \begin{bmatrix} \vec{x}^T & -\vec{x}^T \cdot \text{cam\_pos} \\ \vec{y}^T & -\vec{y}^T \cdot \text{cam\_pos} \\ \vec{z}^T & -\vec{z}^T \cdot \text{cam\_pos} \\ 0\,0\,0 & 1 \end{bmatrix}$ \Comment{orthonormal basis + translation}
\State \textbf{Construct Projection matrix} $P$ (perspective):
\State \quad $P \gets \begin{bmatrix} \frac{1}{\tan(\text{fov}/2) \cdot \text{aspect}} & 0 & 0 & 0 \\ 0 & \frac{1}{\tan(\text{fov}/2)} & 0 & 0 \\ 0 & 0 & \frac{f+n}{n-f} & \frac{2fn}{n-f} \\ 0 & 0 & -1 & 0 \end{bmatrix}$
\For{each joint position $p_i$ (in homogeneous coords)}
    \State $p_{\text{clip}} \gets P \cdot V \cdot p_i$ \Comment{$4\times4$ matrix multiplications}
    \State $p_{\text{screen}} \gets (p_{\text{clip}}.x / p_{\text{clip}}.w,\, p_{\text{clip}}.y / p_{\text{clip}}.w)$ \Comment{perspective divide}
\EndFor
\State Draw line segments between consecutive projected points $\{p_{\text{screen},0}, \dots, p_{\text{screen},n}\}$
\end{algorithmic}
\end{algorithm}

\section{Testing Plan}

We will validate the implementation through the following tests:

\begin{enumerate}[nosep]
    \item \textbf{FK consistency}: set known joint angles (e.g., all zeros), verify the end-effector position matches the expected sum of link lengths along the initial axis.
    \item \textbf{IK convergence}: place the target at reachable positions and check that the error $\|p_{\text{end}} - p_{\text{target}}\|$ decreases monotonically and drops below $\varepsilon$ within a reasonable number of iterations.
    \item \textbf{Unreachable targets}: set the target outside the arm's workspace and verify the arm fully extends toward it without numerical divergence.
    \item \textbf{Singularity robustness}: move the target so the arm passes through a fully stretched (singular) configuration; verify no NaN or explosion in joint angles.
    \item \textbf{Projection correctness}: compare projected 2D coordinates against manually computed values for simple camera setups.
\end{enumerate}

No external datasets are needed. All test data (target positions, arm configurations) is generated programmatically. The visualization itself serves as a continuous qualitative test---unnatural motion or rendering artifacts are immediately visible.

\section{Plan of Future Work}

\begin{table}[H]
\centering
\begin{tabular}{@{}lll@{}}
\toprule
\textbf{Week} & \textbf{Task} & \textbf{Responsible} \\
\midrule
1--2 & Implement FK with rotation matrices and basic 3-link arm & Nazar \\
3 & Implement Jacobian computation and basic pseudoinverse IK & Nazar \\
4 & Add SVD-based damped least squares for stability & Mykola \\
5 & Implement minimal 3D projection and integrate with simple UI & Roman \\
6 & Testing and bug fixes for core functionality & All \\
\bottomrule
\end{tabular}
\end{table}

\section{Potential Challenges}

\begin{itemize}[nosep]
    \item \textbf{Project scope and workload}: implementing FK, IK with SVD, full 3D rendering pipeline, and interactive UI within a 6-week timeline is substantial. We will focus on core functionality first (3-link arm, basic DLS solver, minimal projection), with advanced features such as multi-DOF arms and sophisticated camera controls deferred as stretch goals if time permits.
    \item \textbf{Numerical instability near singularities}: when the arm is fully extended or folded, the Jacobian becomes rank-deficient. The damped pseudoinverse mitigates this, but tuning $\lambda$ is nontrivial---too large slows convergence, too small reintroduces instability.
    \item \textbf{Local minima}: the Jacobian-based method is a local optimizer. The arm may get stuck in configurations that are locally optimal but far from the global solution. We may need to add random restarts or joint-limit handling.
    \item \textbf{Real-time performance}: SVD computation every frame can be expensive for many joints. We may need to limit the number of IK iterations per frame or use the simpler Jacobian transpose as a fallback.
    \item \textbf{3D interaction}: building intuitive mouse-based 3D target control requires careful mapping from 2D screen coordinates back to 3D space (an inverse projection problem in itself).
\end{itemize}

\begin{thebibliography}{9}
\bibitem{buss2004} S.~Buss, ``Introduction to Inverse Kinematics with Jacobian Transpose, Pseudoinverse and Damped Least Squares Methods,'' 2004. Available: \url{https://math.ucsd.edu/~sbuss/ResearchWeb/ikmethods/iksurvey.pdf}
\bibitem{roboticworx} RoboticWorx, ``Inverse Kinematics.'' Available: \url{https://roboticworx.io/blogs/projects/inverse-kinematics}
\bibitem{mathworks} MathWorks, ``Modeling Inverse Kinematics in a Robotic Arm.'' Available: \url{https://se.mathworks.com/help/fuzzy/modeling-inverse-kinematics-in-a-robotic-arm.html}
\end{thebibliography}

\end{document}
